\chapter{Literature Review}
\label{ch:literature}

\section{Scope and search strategy}

This review is not a survey of machine learning in civil engineering. It is
the evidence base for a specific set of design decisions, and each section
closes by naming the decision it supports or the decision it fails to
support. Forty-three sources are cited, each verified against the
publisher record or a primary index.

Searches were run in August 2026 across Google Scholar, arXiv, the ACM
Digital Library, ScienceDirect, Springer, IEEE Xplore and Google Patents,
using the terms \emph{floor plan dataset}, \emph{floorplan
vectorisation}, \emph{structural layout generation}, \emph{column
placement}, \emph{shear wall layout}, \emph{structural surrogate model},
\emph{topology optimisation deep learning}, \emph{noise contrastive
estimation}, \emph{learning to rank} and \emph{variational quantum
classifier}. Backward snowballing proceeded from two reviews
\citep{sun2021mlreview,malaga2022opinionated} and from the closest prior
art \citep{wang2024columnrl,liao2021shearwallgan}. A source was included
only if it supplies data this work consumes, establishes a method it uses,
or contests a decision it makes.

Three limits on the search should be stated at the outset. Coverage is
English-language and mainly open-access, which under-represents generative
structural design as practised and published in China. Patents were
searched only through Google Patents and only in English, which matters
more than usual here because the closest prior art proved to be a patent.
And this is a narrative review assembled around an existing design rather
than a protocol-driven systematic review, so it is exposed to the bias of
having searched for support; Section~\ref{sec:priorart} is the partial
corrective, since it reports the work that anticipates part of the
approach.

\section{Floor plans as a source of structural data}

Large vector floor-plan corpora are now openly available, and three are
used here. CubiCasa5K \citep{kalervo2019cubicasa} contributes 5,000
densely annotated plans as SVG polygons. Swiss Dwellings
\citep{standfest2022swiss} contributes roughly 42,500 apartment models with
walls stored as geo-referenced geometry, released under CC-BY-4.0 and
quality assured to a median area deviation of 1.2 per cent. ResPlan
\citep{abouagour2025resplan} contributes 17,000 residential plans in
vector form with metric coordinates.

Two further corpora were considered and rejected. RPLAN
\citep{wu2019rplan} is the most cited residential corpus but is
raster-first and oriented towards room semantics, so metric wall geometry
would have to be recovered before anything structural could be said. MSD
\citep{vanengelenburg2024msd} is derived from Swiss Dwellings and would
have duplicated a source already present.

Recovering geometry from drawings is itself a mature field. Liu and
colleagues \citep{liu2017raster2vector} framed raster-to-vector conversion
as junction prediction followed by integer programming, and Zeng and
colleagues \citep{zeng2019deepfloorplan} improved wall and room
recognition with a multi-task network. This work deliberately does not
compete there: all three sources are consumed as vectors, so the parsing
problem begins one step later, at walls to \emph{structural} grid rather
than pixels to walls.

\textbf{What this justifies.} The choice of three sources and the decision
to work from vector geometry. \textbf{What it does not justify} is
treating those corpora as structural ground truth, which
Section~\ref{sec:labelgap} addresses.

\section{Learned structural layout design}

Three strands of work bear on this thesis, and they are not equally close
to it.

\subsection{Response prediction}

The largest strand predicts structural behaviour from a design that
already exists, usually by training a surrogate model on simulated
results. Sun, Burton and Huang \citep{sun2021mlreview} survey the field
and show that prediction and assessment dominate while proposing a layout
is thinly covered. This thesis does not sit in that strand.

M\'alaga-Chuquitaype \citep{malaga2022opinionated} is more pointed,
arguing that the field has been readier to automate the parts of design
engineers do not mind losing than the judgement they exercise, and warning
against models asked to certify what they cannot explain. That warning
shapes the division of labour adopted here: the learned components rank
candidates and nothing more, and the reasoning that orders them is set out
in closed form.

\subsection{Generative structural design}

Liao and colleagues \citep{liao2021shearwallgan} trained generative
adversarial networks on real shear-wall design documents, parameterised by
height and seismic category, and produced usable layouts. Lu and
colleagues \citep{lu2022physicsgan} then embedded a physics estimator in
the generator rather than trusting the data-driven discriminator alone.
This strand is the strongest available evidence that structural layout
\emph{is} learnable from real design data.

Two limits matter here. The training documents are proprietary design
archives rather than open data, so the work cannot be reproduced outside
the group that holds them; and the task is shear-wall layout in Chinese
residential practice rather than column-grid inference from an arbitrary
architectural plan.

\subsection{Column placement, and the closest prior art}
\label{sec:priorart}

The closest prior art is not a paper but a patent. Wang and Nourbakhsh at
Autodesk \citep{wang2024columnrl}, filed October 2020 and granted March
2024, claim a two-agent reinforcement learning system in which the first
agent places grid lines so as to maximise overlap with load-bearing walls
subject to minimum and maximum span, and the second places columns at grid
crossings under rules marking locations as required, preferred, candidate
or forbidden.

That is the same decomposition used here: walls vote for grid lines, span
limits bound the search, columns stand at crossings, and some crossings
are excluded. The patent was found after the pipeline had been built and is
reported rather than omitted. It is strong external evidence that the
decomposition is the natural one, and it sharpens what remains to
contribute, as Table~\ref{tab:patent} sets out.

\begin{table}[H]
\centering
\caption[Comparison with the closest prior art]{Comparison with the
closest prior art, US Patent 11,941,327 B2. The decomposition is conceded;
the contribution lies in the ranking and its evaluation.}
\label{tab:patent}
\small
\begin{tabularx}{\textwidth}{lXX}
\toprule
 & \textbf{US 11,941,327 B2} & \textbf{This work} \\
\midrule
Search & Two reinforcement learning agents & Enumeration over interior
line subsets, with seeded sampling past a stated limit \\
Ranking & Reward terms defined by the authors & Four rankers compared
under one protocol, with the algorithm-only path retained as a baseline \\
Data & Not disclosed & Two open corpora with derived labels, regenerable
end to end, coordinates published \\
Evaluation & Viability rules & Coverage reported first, then paired
significance tests on held-out plans \\
Reproducibility & Patent; no released code or data & Open notebooks and
committed data tables \\
\bottomrule
\end{tabularx}
\end{table}

Their treatment of ``required'' locations also justifies a decision taken
in the parser: columns that the drawing itself marks are input that the
generator must honour, not part of the answer to be predicted.

\subsection{Column layouts from plans with neural networks}

Ampanavos, Nourbakhsh and Cheng \citep{ampanavos2021approxiframer}
published the closest work to this task: a convolutional network that
generates column positions from plan sketches, trained on roughly 137,000
synthetic rectangular frames and evaluated by positional error against
synthetic labels. Their work removes any claim that predicting columns
from plans with a neural network is new. It differs from this thesis on
the two points that matter here: its labels are synthetic rather than
derived from real buildings, and it generates a layout rather than ranking
candidates.

Pizarro and colleagues \citep{pizarro2021walls} predicted shear-wall
layouts from architectural plans using convolutional and conditional
generative models trained on 165 real Chilean projects. This is the
nearest precedent for using real buildings as supervision, and it removes
any claim that doing so is new. Its data is private, its subject is wall
panels rather than columns, and it generates rather than ranks.

\section{Topology optimisation, and why it is the wrong frame}

Deep learning has been applied extensively to accelerate topology
optimisation. Sosnovik and Oseledets \citep{sosnovik2019nn4topopt} treated
the optimisation iteration as image-to-image translation; Yu and
colleagues \citep{yu2019neartopo} removed the iteration altogether and
predicted a near-optimal topology directly; and a recent review
\citep{topopt2023review} maps the field.

This is the closest neighbouring literature and it is the wrong tool for
the present problem. Topology optimisation returns a density field over a
continuum. A column grid is a discrete, dimensionally coordinated object
that must be built, detailed and connected, and the operation that turns a
density field into a buildable grid is precisely the difficult part, which
the density formulation does not address. Enumerating discrete
arrangements under span constraints keeps every candidate buildable by
construction.

\textbf{What this justifies} is the discrete generator of
Section~\ref{sec:generator}. \textbf{What it costs} is scalability, which
Section~\ref{sec:gap} returns to.

\section{Screening cheaply and confirming expensively}

The shape of this pipeline, in which many candidates are screened by a
cheap model and few are examined closely, is not invented here. It is the
standard structure of surrogate-assisted optimisation, surveyed by Jin
\citep{jin2011surrogate}: use a cheap model to filter and spend the real
evaluation budget on the survivors.

What this thesis takes from that literature is the architecture, together
with one deliberate departure from it. In surrogate-assisted optimisation
the learned model sits \emph{inside} the search, standing in for an
expensive evaluation. Here the algorithmic score and the learned judge are
two \emph{independent} rankers of the same candidate set, combined by a
blend whose weight is ablated, and compared against the score-only
baseline. No published head-to-head comparison of that kind was found for
building structural layout, which is why the contribution sits in the
evaluation rather than in the search.

Danhaive and Mueller \citep{danhaive2021subspace} supply the complementary
argument: optimisation alone systematically misses the qualitative
properties designers care about, which is why performance-conditioned
exploration outperforms pure search. That is the case for the judges
existing at all alongside a closed-form score.

\section{How the learned components are set up}

Four choices in the judges and in the learned score weights follow
directly from established method.

\textbf{Learning from constructed negative examples.} No dataset labels a
layout as implausible, so the judges construct their negatives: real
arrangements are the positive examples and altered copies of those same
arrangements are the negatives. The method is noise contrastive estimation
\citep{gutmann2010nce}, and naming it matters because it carries a
warning: what the model learns is the boundary between the real data and
\emph{the particular negatives that were constructed}. How the alteration
is performed is therefore a design decision rather than an implementation
detail, and Section~\ref{sec:negatives} shows what changed when that was
taken seriously.

\textbf{Ranking within a group.} The learned score weights are not fitted
by asking whether an arrangement is real in isolation, but by asking which
of several arrangements on one plan is the one that plan uses. That is
learning to rank \citep{liu2009ltr}, in which a set of candidates belongs
to a single query; here the plan is the query.

\textbf{Splitting the data so the reported score is honest.} If rows are
split at random, a plan's real arrangement can fall in the training set
while an altered copy of it falls in the test set; the model then
recognises the copy and scores well without having learned anything
general. Ecologists encountered the same problem with spatially clustered
samples, and Roberts and colleagues \citep{roberts2017cv} showed that
random splits inflate scores and that data should be split by whatever the
samples have in common. Kaufman and colleagues \citep{kaufman2012leakage}
give the general name, leakage. Every model here is split by plan and the
split is repeated over several seeds.

\textbf{Choice of model.} Gradient boosting \citep{friedman2001gbm}, in
its fast histogram form \citep{ke2017lightgbm}. On tables of a few
thousand rows, tree ensembles still outperform neural networks
\citep{grinsztajn2022tabular}, so declining to use deep learning for the
tabular judge is an evidence-based decision rather than an omission. The
implementation is scikit-learn \citep{pedregosa2011sklearn}.

\section{What the industry cost evidence still justifies}

Section~\ref{sec:removals} explained why the pricing was removed. The
evidence that motivated it remains useful at a higher level of
abstraction, because what it establishes is which geometric quantities
drive cost \citep{aisc_cost}: that making and erecting steel is about
seventy per cent of a steel package, and that connections drive thirty to
fifty per cent of fabrication cost while contributing under five per cent
of the weight. Member count, which scales with column count, is therefore
a first-order driver, and columns standing away from walls incur
disproportionate detailing effort. ISO 2848 \citep{iso2848} makes modular
coordination a published standard, so repeated bay sizes are a real habit
of real buildings and a real economy. UK guidance \citep{sci_grids}
supplies the economical span band for offices, which sets the default span
limits in the demonstration application.

Each of those becomes one measure of the score in
Section~\ref{sec:score}: the drivers are retained and the currency is
discarded.

\section{Quantum machine learning, and the limits of what can be claimed}

\textbf{The model type is established.} Havl\'i\v{c}ek and colleagues
\citep{havlicek2019} introduced both the variational classifier and the
quantum kernel and ran them on hardware. Schuld and colleagues
\citep{schuld2020circuit} defined the circuit-centric design this work
starts from, and Schuld and Killoran
\citep{schuld2019featurehilbert} explain why encoding a number as a
rotation angle constitutes a feature map, which is why inputs are scaled
into a bounded angle range. PennyLane \citep{bergholm2018pennylane} is the
software. P\'erez-Salinas and colleagues
\citep{perezsalinas2020reupload} showed that feeding the data into the
circuit again at every layer makes a small circuit considerably more
expressive, and both architectures are tested here.

\textbf{The limits are equally established and set the size of the
circuit.} McClean and colleagues \citep{mcclean2018barren} showed that for
many circuit designs the training signal vanishes exponentially as qubits
are added, which is why this work stays small. Huang and colleagues
\citep{huang2021powerdata} showed that classical models with data are
competitive even on problems constructed to favour quantum ones. Bowles,
Ahmed and Schuld \citep{bowles2024benchmark} tested twelve quantum models
across 160 datasets, found classical models generally ahead, and set the
standard of fairness this work follows: identical rows, identical
features, identical splits, with classical baselines run on the same data.

\textbf{The nearest applications are not design.} Quantum machine learning
has reached structural engineering as post-earthquake safety
classification. Bhatta and Dang \citep{bhatta2024quantum} train a
variational quantum classifier on tabular building and demand features and
compare it against ten classical models on identical data. That is the
closest quantum work to this thesis, and the reason the novelty claim in
Section~\ref{sec:contribution} is worded as design-stage arrangement
plausibility rather than as quantum learning applied to buildings.
Elsewhere, quantum treatment of structural geometry exists as
annealing-based truss optimisation and as eigenvalue analysis, and a
review of the field \citep{ploennigs2026quantumcivil} records no
application to design-stage layout.

\section{Judging the same layout two ways}

Comparisons of representation exist for \emph{generative} models, and
comparisons of tabular models against convolutional networks exist for
floor plans in the prediction of rent rather than plausibility. Van
Engelenburg and colleagues \citep{vanengelenburg2023ssig} propose a
non-learned metric for floor-plan similarity and benchmark it against deep
metric learning, which is the nearest heuristic-versus-learned comparison
on this kind of data, though the task is retrieval rather than
plausibility.

Evaluation of generated structural layouts exists in several forms. Wang
and colleagues \citep{wang2025experteval} represent the state of the art
in selecting among candidate structural layouts: an expert-weighted
evaluation of automatically generated shear-wall schemes combined with
automated finite-element analysis. Its evaluator applies expert-assigned
weights rather than a model learned from data, and no learned judge is
compared against that evaluator on the same candidates. Outside the
construction domain, the recipe of real examples as positives and
perturbed examples as negatives is established for judging object
placement in images \citep{liu2024diffpop}.

\section{The gap}
\label{sec:gap}
\label{sec:labelgap}

The general form of the claim is false: it is not true that structural
layout has never been learned from real data, since
\citet{liao2021shearwallgan} and \citet{pizarro2021walls} have done so.
The claim that survives the evidence is narrower and has four parts.

\begin{enumerate}
\item \textbf{Open architectural collections have not been used as
structural training data.} They are large, free and reproducible, and they
are used for room-layout research. Learned structural design uses private
archives or geometry generated for the purpose.

\item \textbf{The problem of obtaining column-layout labels has not been
solved publicly.} Architectural drawings do not record a column layout, so
anyone learning columns from plans must generate the answer synthetically
\citep{ampanavos2021approxiframer}, hold a private engineering archive
\citep{pizarro2021walls}, or derive a label. No derivation, with its
assumptions and their sensitivity, has been published for the open
corpora. Chapter~\ref{ch:method} is this work's answer, and its 160\,mm
threshold is varied rather than asserted.

\item \textbf{Candidate rankers are not compared with one another.}
Generated or enumerated layouts are evaluated by rules, by finite-element
analysis, or by expert weighting \citep{wang2025experteval}, but no study
was found that ranks one fixed candidate set with both an algorithmic
score and independent learned judges and reports whether the judges
improve on the score alone. The four-variant protocol is this work's
answer.

\item \textbf{Exhaustive checking does not scale, and no remedy is
established.} Exhaustive search is exact on simple plans and infeasible on
heavily subdivided ones. Sampling beyond a stated limit is a mitigation
rather than a solution, and the generative literature scales but does not
produce buildable grids.
\end{enumerate}

Items 1, 2 and 3 are where this thesis contributes. Item 4 remains open
and is presented as such.

One further caveat belongs with the novelty claim. It rests in part on
absence of evidence: searches across several databases found no comparable
study, but no search establishes that none exists. The claim is therefore
worded so that a single counter-example would falsify it, and the works
named in this chapter are those that would come closest to providing one.
